\documentclass[11pt]{article}
\usepackage[margin=1in]{geometry}
\usepackage[T1]{fontenc}
\usepackage[utf8]{inputenc}

% Core packages
\usepackage{amsmath,amssymb}
\usepackage{array,longtable,booktabs}
\usepackage{xcolor}
\usepackage{listings}
\usepackage{hyperref}

% Define colors for code
\definecolor{codegreen}{rgb}{0,0.6,0}
\definecolor{codegray}{rgb}{0.5,0.5,0.5}
\definecolor{codepurple}{rgb}{0.58,0,0.82}
\definecolor{backcolour}{rgb}{0.95,0.95,0.92}

% Python style
\lstdefinestyle{pythonstyle}{
    language=Python,
    backgroundcolor=\color{backcolour},   
    commentstyle=\color{codegreen},
    keywordstyle=\color{magenta},
    numberstyle=\tiny\color{codegray},
    stringstyle=\color{codepurple},
    basicstyle=\ttfamily\footnotesize,
    breakatwhitespace=false,         
    breaklines=true,                 
    captionpos=b,                    
    keepspaces=true,                 
    numbers=left,                    
    numbersep=5pt,                  
    showspaces=false,                
    showstringspaces=false,
    showtabs=false,                  
    tabsize=4
}

% Bash style
\lstdefinestyle{bashstyle}{
    language=bash,
    backgroundcolor=\color{backcolour},   
    commentstyle=\color{codegreen},
    keywordstyle=\color{blue},
    numberstyle=\tiny\color{codegray},
    stringstyle=\color{codepurple},
    basicstyle=\ttfamily\footnotesize,
    breakatwhitespace=false,         
    breaklines=true,                 
    captionpos=b,                    
    keepspaces=true,                 
    numbers=left,                    
    numbersep=5pt,                  
    showspaces=false,                
    showstringspaces=false,
    showtabs=false,                  
    tabsize=4
}

% Plain text style for Markdown
\lstdefinestyle{textstyle}{
    backgroundcolor=\color{backcolour},   
    basicstyle=\ttfamily\footnotesize,
    breakatwhitespace=false,         
    breaklines=true,                 
    captionpos=b,                    
    keepspaces=true,                 
    numbers=left,                    
    numbersep=5pt,                  
    showspaces=false,                
    showstringspaces=false,
    showtabs=false,                  
    tabsize=4
}

\lstset{
    literate=
        {—}{{\textemdash}}1
        {🛑}{{[STOP]}}1
        {✅}{{[OK]}}1
}

\title{Kurdish Medical Dataset Generation with Ollama}
\author{Research Implementation}
\date{\today}

\begin{document}

\maketitle

\tableofcontents
\newpage

\section{Original Request and System Information (prompt.md)}
{\small
\renewcommand{\arraystretch}{1.2}
\begin{longtable}{@{}>{\raggedright\arraybackslash}p{0.28\textwidth}>{\raggedright\arraybackslash}p{0.64\textwidth}@{}}
\caption{Structured summary of \texttt{prompt.md}} \\
\toprule
\textbf{Item} & \textbf{Details} \\
\midrule
\endfirsthead
\toprule
\textbf{Item} & \textbf{Details} \\
\midrule
\endhead
\bottomrule
\endfoot
\bottomrule
\endlastfoot
\multicolumn{2}{@{}l}{\textbf{Original Request}} \\
\midrule
Requested change & change this code from llamacpp to ollama for installation and running start the server and even file name \\
Testing constraint & pleases don't run any code for test in my machine because i copy this file to my server just write the code and i will run it in my server this information about my server in below \\
Model note & also i alreaded downloaded gemma 3 27 b \\
\addlinespace
\multicolumn{2}{@{}l}{\textbf{Operating System \& Kernel}} \\
\midrule
System & Linux \\
Node Name & a214-Z790-D \\
Release & 5.15.0-139-generic \\
Version & \#149\textasciitilde{}20.04.1-Ubuntu SMP Wed Apr 16 08:29:56 UTC 2025 \\
Machine & x86\_64 \\
Processor & x86\_64 \\
\addlinespace
\multicolumn{2}{@{}l}{\textbf{CPU Information}} \\
\midrule
Physical cores & 14 \\
Total cores & 20 \\
Max Frequency & 4620.00~MHz \\
Current Frequency & 3257.71~MHz \\
CPU Usage & 4.4\% \\
\addlinespace
\multicolumn{2}{@{}l}{\textbf{GPU Information}} \\
\midrule
GPU 0 & NVIDIA GeForce RTX 4090 D, Load: 0.0\%, Free Mem: 48123.0~MB, Used Mem: 405.0~MB, Total Mem: 49140.0~MB, Temperature: 32.0~$^\circ$C \\
GPU 1 & NVIDIA GeForce RTX 4090 D, Load: 0.0\%, Free Mem: 48511.0~MB, Used Mem: 25.0~MB, Total Mem: 49140.0~MB, Temperature: 29.0~$^\circ$C \\
\addlinespace
\multicolumn{2}{@{}l}{\textbf{Memory Information}} \\
\midrule
Total & 15.44~GB \\
Available & 11.40~GB \\
Used & 4.04~GB \\
Percentage & 26.2\% \\
\end{longtable}
}

\newpage
\section{Installation Scripts}

\subsection{install\_ollama.sh}
\begin{lstlisting}[style=bashstyle, caption={install\_ollama.sh}]
#!/bin/bash
# install_ollama.sh - Install Ollama on Ubuntu (CUDA/GPU enabled)
set -e

echo "================================================"
echo "  Ollama Installation Script"
echo "  Target: Ubuntu 20.04 | Dual RTX 4090 D"
echo "================================================"

# 1. Install CURL if not present
echo "[*] Ensuring curl is installed..."
sudo apt-get update -y
sudo apt-get install -y curl

# 2. Remove any existing Ollama installation before reinstalling
echo ""
echo "[*] Checking for existing Ollama installation..."
OLLAMA_BIN=$(command -v ollama || echo "/usr/local/bin/ollama")

if [ -f "$OLLAMA_BIN" ] || systemctl list-unit-files | grep -q ollama.service; then
    echo "[!] Ollama installation detected — removing it completely for a clean reinstall..."

    # Stop and disable the service
    sudo systemctl stop ollama 2>/dev/null || true
    sudo systemctl disable ollama 2>/dev/null || true

    # Remove the binary
    if [ -f "$OLLAMA_BIN" ]; then
        sudo rm -f "$OLLAMA_BIN"
        echo "    [v] Binary removed ($OLLAMA_BIN)"
    fi
    # Also check /usr/bin just in case
    sudo rm -f /usr/bin/ollama

    # Remove the systemd unit and any drop-in overrides
    sudo rm -f /etc/systemd/system/ollama.service
    sudo rm -rf /etc/systemd/system/ollama.service.d
    sudo systemctl daemon-reload
    echo "    [v] Systemd unit and overrides removed"

    # Remove system data directory (libraries/runners)
    if [ -d /usr/share/ollama ]; then
        sudo rm -rf /usr/share/ollama
        echo "    [v] System data directory /usr/share/ollama removed"
    fi

    # Remove the ollama system user/group
    if id -u ollama &>/dev/null; then
        sudo userdel ollama 2>/dev/null || true
        echo "    [v] System user 'ollama' removed"
    fi
    if getent group ollama &>/dev/null; then
        sudo groupdel ollama 2>/dev/null || true
        echo "    [v] System group 'ollama' removed"
    fi

    echo "[+] Old Ollama installation removed."
else
    echo "[+] No existing Ollama installation found — proceeding with fresh install."
fi
echo ""

# 3. Fresh Ollama install (automatically detects CUDA and installs GPU support)
echo "[*] Running official Ollama installer..."
curl -fsSL https://ollama.com/install.sh | sh

# 3. Configure Ollama storage on the large drive (/home is full)
echo "[*] Creating Ollama storage directory on /mnt/storage1/shko/ollama ..."
export OLLAMA_HOME="/mnt/storage1/shko/ollama"
sudo mkdir -p "$OLLAMA_HOME"
# Make the directory owned by the user that will run ollama
sudo chown -R "$USER":"$USER" "$OLLAMA_HOME"

# 4. Configure Ollama to listen on all interfaces (so the Python client can reach it)
echo "[*] Configuring Ollama systemd service (host, GPUs, storage)..."
sudo mkdir -p /etc/systemd/system/ollama.service.d
cat <<EOF | sudo tee /etc/systemd/system/ollama.service.d/override.conf
[Service]
Environment="OLLAMA_HOST=0.0.0.0:11434"
# Use both GPUs
Environment="CUDA_VISIBLE_DEVICES=0,1"
# Store all Ollama data (models, blobs, keys) on the large drive
Environment="OLLAMA_HOME=/mnt/storage1/shko/ollama"
EOF

# 5. Reload systemd and start the Ollama service
echo "[*] Enabling and starting Ollama service..."
sudo systemctl daemon-reload
sudo systemctl enable ollama
sudo systemctl restart ollama

echo ""
echo "================================================"
echo "  SUCCESS! Ollama is installed and running."
echo "  Service status: sudo systemctl status ollama"
echo "  API endpoint  : http://localhost:11434"
echo "================================================"
echo ""
echo "  NEXT STEPS:"
echo "  1. Run ./setup_ollama_model.sh    <- registers Gemma 3 27B"
echo "  2. Run python3 run_server_bg_first.py  <- starts the generator"
echo "================================================"
\end{lstlisting}

\newpage
\subsection{setup\_ollama\_model.sh}
\begin{lstlisting}[style=bashstyle, caption={setup\_ollama\_model.sh}]
#!/bin/bash
# setup_ollama_model.sh - Register the local Gemma 3 27B GGUF with Ollama
# Run this AFTER install_ollama.sh
#
# The model weight file is expected at:
#   model_weights/gemma-3-27b-it-q4_k_m.gguf
# (relative to THIS script's directory, i.e. the directory you uploaded to the server)
#
# Ollama's "create" command reads a Modelfile, imports the GGUF blob into its
# internal content-addressed store, and registers it as "gemma3:27b".
# After this runs, "ollama run gemma3:27b" works offline — no internet needed.

set -e

# -- Storage --------------------------------------------------------------------
export OLLAMA_HOME="/mnt/storage1/shko/ollama"
export OLLAMA_MODELS="/mnt/storage1/shko/ollama/Models"
mkdir -p "$OLLAMA_HOME"
mkdir -p "$OLLAMA_MODELS"

# -- Paths ----------------------------------------------------------------------
SCRIPT_DIR="$(cd "$(dirname "${BASH_SOURCE[0]}")" && pwd)"
GGUF_FILE="$SCRIPT_DIR/model_weights/gemma-3-27b-it-q4_k_m.gguf"
MODEL_NAME="gemma3:27b"
MODELFILE="$SCRIPT_DIR/Modelfile"

echo "================================================"
echo "  Ollama Model Setup (local GGUF import)"
echo "  Model  : $MODEL_NAME"
echo "  GGUF   : $GGUF_FILE"
echo "  Home   : $OLLAMA_HOME"
echo "  Models : $OLLAMA_MODELS"
echo "================================================"

# -- Sanity checks --------------------------------------------------------------
if [ ! -f "$GGUF_FILE" ]; then
    echo ""
    echo "[!] ERROR: GGUF file not found at:"
    echo "    $GGUF_FILE"
    echo ""
    echo "    Make sure you transferred the file with the correct path:"
    echo "    model_weights/gemma-3-27b-it-q4_k_m.gguf"
    exit 1
fi

if ! curl -s http://localhost:11434/api/tags > /dev/null 2>&1; then
    echo "[!] Ollama service does not seem to be running."
    echo "    Start it with:  sudo systemctl start ollama"
    echo "    Or manually  :  OLLAMA_HOME=$OLLAMA_HOME ollama serve"
    exit 1
fi

# -- Write a minimal Modelfile --------------------------------------------------
echo "[*] Writing Modelfile..."
cat > "$MODELFILE" <<EOF
# Modelfile for Gemma 3 27B Instruct (Q4_K_M GGUF)
FROM $GGUF_FILE

# Recommended chat template for Gemma 3 instruct models
TEMPLATE """<start_of_turn>user
{{ .Prompt }}<end_of_turn>
<start_of_turn>model
"""

# Generation parameters
PARAMETER temperature    0.7
PARAMETER top_p          0.9
PARAMETER top_k          40
PARAMETER repeat_penalty 1.1
PARAMETER num_ctx        8192
EOF

echo "    Modelfile written to: $MODELFILE"

# -- Import the model -----------------------------------------------------------
echo ""
echo "[*] Running: ollama create $MODEL_NAME -f $MODELFILE"
echo "    This copies the GGUF into Ollama's blob store — may take a minute..."
echo ""

OLLAMA_HOME="$OLLAMA_HOME" OLLAMA_MODELS="$OLLAMA_MODELS" ollama create "$MODEL_NAME" -f "$MODELFILE"

echo ""
echo "================================================"
echo "  SUCCESS! Model is ready."
echo ""
echo "  Quick test:"
echo "    OLLAMA_HOME=$OLLAMA_HOME OLLAMA_MODELS=$OLLAMA_MODELS ollama run $MODEL_NAME \"Hello!\""
echo ""
echo "  List registered models:"
echo "    OLLAMA_HOME=$OLLAMA_HOME OLLAMA_MODELS=$OLLAMA_MODELS ollama list"
echo ""
echo "  The model is now registered as: $MODEL_NAME"
echo "  Start the server with: python3 run_server_bg_first.py"
echo "================================================"
\end{lstlisting}

\newpage
\subsection{install\_requirements.sh}
\begin{lstlisting}[style=bashstyle, caption={install\_requirements.sh}]
#!/bin/bash
# install_requirements.sh - Python requirements for dataset generation
# Ollama runs as its own service, so we only need the OpenAI-compatible client + helpers.

set -e

echo "Updating system packages..."
sudo apt-get update

echo "Installing system dependencies..."
# psmisc provides 'fuser', used in the background manager scripts
sudo apt-get install -y psmisc python3-pip python3-venv

echo "Installing Python requirements for the project..."
# openai  - speaks to Ollama's /v1 OpenAI-compatible endpoint
# httpx   - used for fine-grained connection control in the generator
pip3 install openai httpx

echo "------------------------------------------------"
echo "Project requirements installed successfully!"
echo ""
echo "NOTE: Ollama itself is the model server."
echo "      No llama-cpp-python, huggingface_hub, or"
echo "      GGUF management libraries are needed."
echo "------------------------------------------------"
\end{lstlisting}

\newpage
\section{Server Management}

\subsection{start\_ollama\_server\_first.sh}
\begin{lstlisting}[style=bashstyle, caption={start\_ollama\_server\_first.sh}]
#!/bin/bash
# start_ollama_server_first.sh
# Ensures the Ollama daemon is running and the Gemma 3 27B model is loaded.
# Ollama itself is the server — this script just makes sure it is up.

# Bypass system proxy for localhost connections
export no_proxy="localhost,127.0.0.1,0.0.0.0"
export NO_PROXY="localhost,127.0.0.1,0.0.0.0"

# -- Storage: redirect all Ollama data to the large drive ----------------------
# /home is full; /mnt/storage1/shko/ollama has the free space.
# OLLAMA_HOME  controls where Ollama stores its config and ed25519 key.
# OLLAMA_MODELS controls where model blobs (the actual weights) are stored.
export OLLAMA_HOME="/mnt/storage1/shko/ollama"
export OLLAMA_MODELS="/mnt/storage1/shko/ollama/Models"
mkdir -p "$OLLAMA_HOME"
mkdir -p "$OLLAMA_MODELS"

# Context window: allow up to 32 K tokens for both input and output
export OLLAMA_NUM_CTX=32768

# Configuration
MODEL_NAME="gemma3:27b"
OLLAMA_HOST="0.0.0.0"
OLLAMA_PORT=11434

echo "------------------------------------------------"
echo "  Starting / checking Ollama server"
echo "  Model : $MODEL_NAME"
echo "  Host  : $OLLAMA_HOST"
echo "  Port  : $OLLAMA_PORT"
echo "  GPUs  : Both RTX 4090 D (CUDA_VISIBLE_DEVICES=0,1)"
echo "------------------------------------------------"

# Export so the Ollama process uses both GPUs
export CUDA_VISIBLE_DEVICES=0,1
export OLLAMA_HOST="${OLLAMA_HOST}:${OLLAMA_PORT}"

# If Ollama is running as a systemd service, restart it to pick up env vars
if systemctl is-active --quiet ollama 2>/dev/null; then
    echo "[*] Ollama systemd service is active — restarting to apply env vars..."
    # Patch the systemd override so the service also knows about OLLAMA_HOME
    sudo mkdir -p /etc/systemd/system/ollama.service.d
    sudo tee /etc/systemd/system/ollama.service.d/storage.conf > /dev/null <<EOF
[Service]
Environment="OLLAMA_HOME=/mnt/storage1/shko/ollama"
Environment="OLLAMA_MODELS=/mnt/storage1/shko/ollama/Models"
EOF
    sudo systemctl daemon-reload
    sudo systemctl restart ollama
    sleep 3
else
    # Otherwise start it as a foreground process (useful when running manually)
    echo "[*] Starting Ollama serve in the foreground..."
    exec ollama serve
fi

# Wait until the REST API is reachable
echo "[*] Waiting for Ollama API to become available..."
for i in $(seq 1 30); do
    if curl -s http://localhost:${OLLAMA_PORT}/api/tags > /dev/null 2>&1; then
        echo "[+] Ollama API is up!"
        break
    fi
    echo "    ... attempt $i/30"
    sleep 2
done

# Pre-load the model so the first inference call is instant
echo "[*] Pre-loading model $MODEL_NAME into GPU VRAM..."
ollama run "$MODEL_NAME" "" 2>/dev/null || true

echo ""
echo "------------------------------------------------"
echo "  Ollama is ready."
echo "  API endpoint : http://localhost:${OLLAMA_PORT}"
echo "  OpenAI compat: http://localhost:${OLLAMA_PORT}/v1"
echo "  Model        : $MODEL_NAME"
echo "------------------------------------------------"
\end{lstlisting}

\newpage
\subsection{run\_server\_bg\_first.py}
\begin{lstlisting}[style=pythonstyle, caption={run\_server\_bg\_first.py}]
import subprocess
import os
import time
import sys
from pathlib import Path

# --- Configuration ---
BASE_DIR = Path(__file__).parent.resolve()
SERVER_SCRIPT = BASE_DIR / "start_ollama_server_first.sh"
SERVER_LOG    = BASE_DIR / "ollama_server_logs_first.txt"

def start_process(command, log_file_handle, description, cwd=None):
    """Starts a process in the background and redirects output to a log file."""
    print(f"[*] Starting {description}...")
    
    process = subprocess.Popen(
        command,
        stdout=log_file_handle,
        stderr=subprocess.STDOUT,
        preexec_fn=os.setpgrp if sys.platform != "win32" else None,
        cwd=str(cwd) if cwd else str(BASE_DIR)
    )
    
    print(f"[+] {description} started with PID: {process.pid}")
    return process

def main():
    if not SERVER_SCRIPT.exists():
        print(f"ERROR: {SERVER_SCRIPT} not found.")
        sys.exit(1)

    # --- Cleanup: stop any stale ollama processes that are NOT the systemd service ---
    print("[*] Cleaning up any stale non-systemd ollama serve processes...")
    os.system("pkill -9 -f 'ollama serve' > /dev/null 2>&1 || true")
    
    print("[*] Waiting 3 seconds for resources to be released...")
    time.sleep(3)

    os.chmod(SERVER_SCRIPT, 0o755)

    print(f"{'='*60}")
    print(f" Ollama Server Background Manager")
    print(f"{'='*60}")
    print(f" Server Log : {SERVER_LOG}")
    print(f"{'='*60}")

    with open(SERVER_LOG, "a") as server_handle:
        server_handle.write(f"\n\n{'='*80}\n")
        server_handle.write(f" OLLAMA SERVER SESSION STARTED AT: {time.strftime('%Y-%m-%d %H:%M:%S')}\n")
        server_handle.write(f"{'='*80}\n\n")
        server_handle.flush()

        server_proc = start_process(["bash", str(SERVER_SCRIPT)], server_handle, "Ollama Server", cwd=BASE_DIR)

    print(f"\n{'='*60}")
    print(" SUCCESS: Ollama server script is running in the background.")
    print(f" Server PID : {server_proc.pid}")
    print(f" Check logs : tail -f {SERVER_LOG}")
    print(f" API        : http://localhost:11434")
    print(f" OpenAI API : http://localhost:11434/v1")
    print(f"{'='*60}\n")
    print(f" To stop, use: kill {server_proc.pid}")
    print(f"  or:          sudo systemctl stop ollama")

if __name__ == "__main__":
    try:
        main()
    except Exception as e:
        print(f"CRITICAL ERROR: {e}")
        sys.exit(1)
\end{lstlisting}

\newpage
\section{Dataset Generation}

\subsection{generate\_dataset\_first.py}
\begin{lstlisting}[style=pythonstyle, caption={generate\_dataset\_first.py}]
#!/usr/bin/env python3
"""
generate_dataset_first.py (Ollama version) - First Dataset Instance

For every .txt file under the configured DATA_DIR:
  - Standalone files  (e.g. 1-10.txt):   run the prompt NUM_RUNS times -> JSON files
                                           -> merge all runs into one final JSON file
  - Part files        (e.g. 163267_part1.txt, 163267_part2.txt, ...):
                        run the prompt NUM_RUNS times per part -> JSON files per part
                        -> after ALL parts are done, merge everything into one final JSON

Original .txt files are NEVER deleted.
Duplicate (question, response) pairs within the same final merged file are removed.

Requires: Ollama running on localhost:11434 with gemma3:27b loaded.
"""

import os
import re
import json
import time
import httpx
from openai import OpenAI

# --------------------------- Bypass Proxy -------------------------------------
# Ensure localhost connections are NOT routed through any system proxy
os.environ["no_proxy"] = "localhost,127.0.0.1"
os.environ["NO_PROXY"] = "localhost,127.0.0.1"

# --------------------------- Configuration ------------------------------------

# Use absolute paths for input and output
DATA_DIR   = "UKH-AIIC-KA_Just_TXT_file/part1"
OUTPUT_DIR = "UKH-AIIC-KA_Just_TXT_file_Output/part1" 
MODEL_NAME = "gemma3:27b"
NUM_RUNS   = 1

# Ollama OpenAI-compatible endpoint
client = OpenAI(
    base_url="http://localhost:11434/v1",
    api_key="ollama",          # Ollama ignores the key but the field must be non-empty
    http_client=httpx.Client(proxy=None),
)

# ---------------------------- Prompt Template ---------------------------------

def build_prompt(text: str, run_index: int) -> str:
    return f"""Act as an expert Kurdish linguist and data engineer.
Analyze the following text and determine its category (e.g., History, Medical, Mathematics, Technology, Literature, Science, etc.).
Then, generate high-quality, information-rich question-and-answer pairs based on the text.

IMPORTANT: Return a VALID JSON array. Start your response with '[' and end with ']'. 
Do NOT include any markdown formatting, code fences (```json), or introductory/concluding text.

Structure:
[
  {{
    "id": "1",
    "category": "<category>",
    "question": "<formal Sorani question>",
    "response": "<formal Sorani answer>",
    "document": {{
      "title": "<title>",
      "source_url": "",
      "publication_date": ""
    }}
  }}
]

Rules:
- Language: Central Kurdish (Sorani) ONLY.
- Grammar: Formal and academic level.
- Format: Output ONLY the raw JSON array.

Source text:
\"\"\"
{text}
\"\"\"
"""

# ------------------------------- Helpers --------------------------------------

def call_model(prompt: str) -> str:
    retry_wait = 30
    while True:
        try:
            response = client.chat.completions.create(
                model=MODEL_NAME,
                messages=[{"role": "user", "content": prompt}],
                temperature=0.7,
                max_tokens=32768,
            )
            return response.choices[0].message.content
        except (httpx.ConnectError, httpx.ConnectTimeout, httpx.RemoteProtocolError) as e:
            print(f"\n    x  Server connection lost: {e}")
            time.sleep(retry_wait)
        except Exception as e:
            print(f"\n    x  Error calling model: {e}")
            time.sleep(retry_wait)

def extract_json_array(raw: str) -> list:
    if not raw or not isinstance(raw, str):
        return []
    cleaned = re.sub(r"```(?:json)?\s*", "", raw)
    cleaned = cleaned.replace("```", "").strip()
    start = cleaned.find("[")
    end   = cleaned.rfind("]")
    if start != -1 and end != -1 and end > start:
        json_str = cleaned[start : end + 1]
        json_str = re.sub(r",\s*([\]}])", r"\1", json_str)
        try:
            data = json.loads(json_str)
            if isinstance(data, list): return data
            if isinstance(data, dict): return [data]
        except json.JSONDecodeError: pass
    return []

def deduplicate(records: list) -> list:
    seen = set()
    unique = []
    for rec in records:
        key = (rec.get("question", "").strip(), rec.get("response", "").strip())
        if key not in seen and key != ("", ""):
            seen.add(key)
            unique.append(rec)
    return unique

def re_index(records: list, offset: int = 0) -> list:
    for i, rec in enumerate(records, start=offset + 1):
        rec["id"] = str(i)
    return records

def save_json(path: str, data: list):
    with open(path, "w", encoding="utf-8") as f:
        json.dump(data, f, ensure_ascii=False, indent=2)
    print(f"    [SAVE] Saved {len(data)} records -> {os.path.relpath(path)}")

def load_json(path: str) -> list:
    with open(path, "r", encoding="utf-8") as f:
        return json.load(f)

def process_single_txt(txt_path: str, output_dir: str, base_name: str):
    print(f"\n  [FILE] Processing: {os.path.relpath(txt_path)}")
    with open(txt_path, "r", encoding="utf-8") as f:
        text = f.read()
    run_files = []
    for run in range(1, NUM_RUNS + 1):
        run_output = os.path.join(output_dir, f"{base_name}_run{run}.json")
        if os.path.exists(run_output):
            try:
                existing = load_json(run_output)
                if len(existing) > 0:
                    run_files.append(run_output)
                    continue
            except: pass
        prompt  = build_prompt(text, run)
        raw     = call_model(prompt)
        records = extract_json_array(raw)
        if records:
            save_json(run_output, records)
            run_files.append(run_output)
        time.sleep(2)
    all_records = []
    for rf in run_files:
        if os.path.exists(rf):
            try: all_records.extend(load_json(rf))
            except: pass
    return re_index(deduplicate(all_records))

def process_folder(input_folder: str, output_folder: str):
    os.makedirs(output_folder, exist_ok=True)
    all_files = sorted(f for f in os.listdir(input_folder) if f.endswith(".txt"))
    standalone = [f for f in all_files if "_part" not in f or f.startswith("part_")]
    parts       = [f for f in all_files if "_part"     in f and not f.startswith("part_")]
    for fname in standalone:
        txt_path  = os.path.join(input_folder, fname)
        base_name = os.path.splitext(fname)[0]
        final_out = os.path.join(output_folder, f"{base_name}_FINAL.json")
        if os.path.exists(final_out): continue
        merged = process_single_txt(txt_path, output_folder, base_name)
        if merged: save_json(final_out, merged)
    if parts:
        groups = {}
        for fname in parts:
            group_key = re.sub(r"_part\d+\.txt$", "", fname)
            groups.setdefault(group_key, []).append(fname)
        for group_key, part_files in groups.items():
            final_out = os.path.join(output_folder, f"{group_key}_FINAL.json")
            if os.path.exists(final_out): continue
            group_records = []
            for fname in sorted(part_files):
                txt_path  = os.path.join(input_folder, fname)
                base_name = os.path.splitext(fname)[0]
                merged = process_single_txt(txt_path, output_folder, base_name)
                group_records.extend(merged)
            if group_records: save_json(final_out, re_index(deduplicate(group_records)))

def main():
    if not os.path.exists(DATA_DIR): return
    process_folder(DATA_DIR, OUTPUT_DIR)

if __name__ == "__main__":
    main()
\end{lstlisting}

\newpage
\subsection{run\_generator\_bg\_first.py}
\begin{lstlisting}[style=pythonstyle, caption={run\_generator\_bg\_first.py}]
import subprocess
import os
import time
import sys
from pathlib import Path

# --- Configuration ---
BASE_DIR = Path(__file__).parent.resolve()
GENERATOR_SCRIPT = BASE_DIR / "generate_dataset_first.py"
GENERATOR_LOG = BASE_DIR / "generate_dataset_logs_first.txt"
PYTHON_EXEC = sys.executable

def start_process(command, log_file_handle, description, cwd=None):
    """Starts a process in the background and redirects output to a log file."""
    print(f"[*] Starting {description}...")
    
    process = subprocess.Popen(
        command,
        stdout=log_file_handle,
        stderr=subprocess.STDOUT,
        preexec_fn=os.setpgrp if sys.platform != "win32" else None,
        cwd=str(cwd) if cwd else str(BASE_DIR)
    )
    
    print(f"[+] {description} started with PID: {process.pid}")
    return process

def main():
    if not GENERATOR_SCRIPT.exists():
        print(f"ERROR: {GENERATOR_SCRIPT} not found.")
        sys.exit(1)

    # --- Selective Cleanup ---
    my_pid = os.getpid()
    print(f"[*] Cleaning up old generator processes for {GENERATOR_SCRIPT.name}...")
    os.system(f"pgrep -f '{GENERATOR_SCRIPT.name}' | grep -v '^{my_pid}$' | xargs -r kill -9 2>/dev/null || true")
    time.sleep(2) 

    print(f"{'='*60}")
    print(f" Dataset Generator Background Manager")
    print(f"{'='*60}")
    print(f" Generator Log : {GENERATOR_LOG}")
    print(f"{'='*60}")

    with open(GENERATOR_LOG, "a") as generator_handle:
        generator_handle.write(f"\n\n{'='*80}\n")
        generator_handle.write(f" DATASET GENERATOR SESSION STARTED AT: {time.strftime('%Y-%m-%d %H:%M:%S')}\n")
        generator_handle.write(f"{'='*80}\n\n")
        generator_handle.flush()

        generator_proc = start_process([PYTHON_EXEC, str(GENERATOR_SCRIPT)], generator_handle, "Dataset Generator", cwd=BASE_DIR)

    print(f"\n{'='*60}")
    print(" SUCCESS: Generator is running in the background.")
    print(f" Generator PID : {generator_proc.pid}")
    print(f" Check logs    : tail -f {GENERATOR_LOG}")
    print(f"{'='*60}\n")
    print(f" To stop, use: kill {generator_proc.pid}")

if __name__ == "__main__":
    try:
        main()
    except Exception as e:
        print(f"CRITICAL ERROR: {e}")
        sys.exit(1)
\end{lstlisting}

\newpage
\section{Testing and Master Script}

\subsection{test\_ollama.sh}
\begin{lstlisting}[style=bashstyle, caption={test\_ollama.sh}]
#!/bin/bash
# test_ollama.sh
# Tests whether the Ollama server and gemma3:27b model are working correctly.

# -- Config --
MODEL_NAME="gemma3:27b"
OLLAMA_PORT=11434
BASE_URL="http://localhost:${OLLAMA_PORT}"

echo "================================================"
echo "  Ollama Test Suite"
echo "================================================"

# -- Test 1: API reachability --
HTTP_CODE=$(curl -s -o /dev/null -w "%{http_code}" "${BASE_URL}/api/tags")
if [ "$HTTP_CODE" = "200" ]; then
    echo "[PASS] Ollama API is reachable"
else
    echo "[FAIL] Ollama API is NOT reachable"
    exit 1
fi

# -- Test 2: Model is listed --
if curl -s "${BASE_URL}/api/tags" | grep -q "${MODEL_NAME}"; then
    echo "[PASS] Model '$MODEL_NAME' found"
else
    echo "[FAIL] Model '$MODEL_NAME' NOT found"
    exit 1
fi

# -- Test 3: Inference --
RESPONSE=$(curl -s -X POST "${BASE_URL}/api/generate" \
    -H "Content-Type: application/json" \
    -d "{
        \"model\": \"${MODEL_NAME}\",
        \"prompt\": \"Reply with confirm you are working.\",
        \"stream\": false
    }")

if echo "$RESPONSE" | grep -q "response"; then
    echo "[PASS] Inference OK"
else
    echo "[FAIL] Inference failed"
fi
\end{lstlisting}

\newpage
\subsection{run.sh}
\begin{lstlisting}[style=bashstyle, caption={run.sh}]
# 1. Install Ollama + GPU support
chmod +x install_ollama.sh && ./install_ollama.sh

# 2. Register / download the Gemma 3 27B model
chmod +x setup_ollama_model.sh && ./setup_ollama_model.sh

# 3. Install Python dependencies
chmod +x install_requirements.sh && ./install_requirements.sh

# 4. Start the Ollama server in background
python3 run_server_bg_first.py

# 5. Start the dataset generator in background
python3 run_generator_bg_first.py

# To stop everything:
chmod +x pkill_ollama.sh && ./pkill_ollama.sh
\end{lstlisting}

\newpage
\subsection{pkill\_ollama.sh}
\begin{lstlisting}[style=bashstyle, caption={pkill\_ollama.sh}]
#!/bin/bash
# pkill_ollama.sh - Stop all Ollama and dataset-generator processes

echo "🛑 Terminating all Ollama and generator processes..."

# Stop the systemd Ollama service (if running)
if systemctl is-active --quiet ollama 2>/dev/null; then
    echo "  [*] Stopping ollama systemd service..."
    sudo systemctl stop ollama
fi

# Kill any manually started ollama serve processes
pkill -9 -f "ollama serve"    2>/dev/null || true
pkill -9 -f "run_server_bg"   2>/dev/null || true
pkill -9 -f "run_generator_bg" 2>/dev/null || true
pkill -9 -f "generate_dataset" 2>/dev/null || true

# Free the Ollama default port just in case
fuser -k 11434/tcp 2>/dev/null || true

echo "✅ All killed."
\end{lstlisting}

\end{document}
